
    \documentclass{standalone}
    \usepackage[english,french]{babel} 
    \usepackage{tikz}
    \usetikzlibrary{calc}
    \usepackage[T1]{fontenc}
    \usepackage[utf8]{inputenc}         
    \usepackage{pgfplotstable}
    \usepackage{pgfplots}
    \pgfplotsset{compat=1.18}
    \usepackage{siunitx}               
    \usepackage{booktabs} 
    \begin{document}
    \begin{tikzpicture}
    \begin{axis}[
                height=10cm,
                width=15cm,
                tick align=outside,
                tick pos=left,
                ylabel={\% territoire},
                ymin= 0,
                ymax= 25,
                ytick={0,5,10,15,20,25},
                yticklabels={0,5,10,15,20,25},
                xtick={1,2,3,4,5,6,7,8,9,10,11,12,13,14,15,16,17,18,19,20,21,22,23,24,25,26,27,28,29,30,31,32,33,34,35},
                xticklabels={SJBA,
STRO,
MCLM,
LAIR,
SSAC,
VLIM,
SSAU,
VQ-CP,
CUNI,
4-6,
MAIZ,
PLAT,
VANI,
PSFO,
JESU,
STLO,
4-5,
NE-LB,
DU-LS,
VIBO,
CHAT,
VIMO,
5-2,
ORSA,
SILR,
CHMO,
CARO,
STEM,
LRTV,
4-3,
VBEL,
NDL,
AÉRO,
LSC,
STLI},
                xlabel={Quartier},
                xmin=-0.5,
                xmax=35.5,xticklabel style={rotate=90, anchor=east}
            ]\addplot[ybar,draw=none,fill=blue,legend image post style={/pgfplots/ybar legend},bar width= 8pt,bar shift=0pt]
                coordinates{
                    (1,22.492163715233396)(2,11.897748208751725)(3,10.165445581174087)(4,8.573918168376824)(5,8.269459106457216)(6,8.229363883410889)(7,6.795597575735954)(8,6.323439058616214)(9,5.569445419471855)(10,5.559808605157528)(11,4.39465026479656)(12,4.277294270247883)(13,4.078765025340388)(14,4.0088279599055285)(15,3.5770802379066726)(16,2.8832615099673005)(17,2.7566773173961288)(18,2.6439094941622483)(19,2.47934287202898)(20,2.3209165350058574)(21,2.30484832933001)(22,2.136151353681366)(23,2.107139060446892)(24,2.028350874280927)(25,1.9908962337416551)(26,1.7162201827329588)(27,1.176278112901378)(28,1.0412344271062488)(29,1.02760288061623)(30,0.9927778640413812)(31,0.5039903889460531)(32,0.36784199428481956)(33,0.3327296673301082)(34,0.33251999410055827)(35,0.27469528334597054)
                };
            \addlegendentry{Pourcentage du territoire};
        \end{axis}
        \end{tikzpicture}
        \end{document} 